\section{Methodology}
\label{sec:methodology}

\subsection{System overview}

\Cref{fig:pipeline} separates training, label-free explanation and mask-based
evaluation. No expert label is used to train the VAE, fit the GMM or produce the
peak ranking.

\begin{figure*}[t]
  \centering
  \resizebox{0.96\textwidth}{!}{%
  \begin{tikzpicture}[
    node distance=1.15cm,
    box/.style={draw, rectangle, rounded corners=2pt, minimum height=1.35cm,
      minimum width=2.55cm, align=center, font=\small, fill=gray!10},
    arrow/.style={-{Stealth[scale=1.15]}, thick},
    evalarrow/.style={-{Stealth[scale=1.15]}, thick, dashed, draw=orange!80!black}
  ]
    \node[box, fill=blue!10] (spectra) {TIC-normalised\\MSI spectra};
    \node[box, fill=orange!12, right=of spectra] (context) {Central spectrum\\$+$ local context};
    \node[box, fill=green!12, right=of context] (vae) {Matched VAE\\variants};
    \node[box, fill=purple!10, right=of vae] (explain) {Latent GMM\\$+$ Integrated Gradients};
    \node[box, fill=blue!10, right=of explain] (peaks) {Ranked\\$m/z$ peaks};
    \node[box, fill=orange!10, below=1.05cm of peaks] (score) {Spatial peak quality\\mSCF1};
    \node[box, fill=gray!8, left=1.15cm of score] (mask) {Expert\\tissue mask};

    \draw[arrow] (spectra) -- (context);
    \draw[arrow] (context) -- (vae);
    \draw[arrow] (vae) -- (explain);
    \draw[arrow] (explain) -- (peaks);
    \draw[arrow] (peaks) -- (score);
    \draw[evalarrow] (mask) -- (score);
  \end{tikzpicture}%
  }
  \caption{Simplified Spatial-msiPL workflow. Each measured spectrum is
  TIC-normalised and supplied either alone or with a summary of its valid
  eight-position Moore neighbourhood. Matched VAE variants reconstruct the
  central spectrum. After training, a label-free two-component GMM is fitted to
  frozen latent means and its posterior is attributed to the original $m/z$
  bins using Integrated Gradients. The resulting ranking is evaluated at a
  matched peak budget. The dashed path emphasises that expert tissue masks are
  used only to calculate mSCF1 after peak selection, never to train the VAE,
  fit the GMM or generate the ranking.}
  \label{fig:pipeline}
\end{figure*}

\subsection{Datasets and preprocessing}

Two public MSI collections were used (\Cref{tab:data}). The MassNet GBM data
contain eight sections with normal/tumour annotations and a shared 85,062-bin
axis~\citep{abdelmoula2022massnet}. Grid coverage ranges from 54.62\% to
82.95\%, so neighbourhoods are constructed from measured coordinates rather
than by treating missing grid locations as spectra. The CAC collection contains
eight sections, 1,481 bins and three mask classes. Seven grids are complete;
\texttt{280TopL} has 4,160 measured pixels in a 4,623-position grid (89.98\%
coverage). Coordinate--mask alignment and HDF5 round trips were audited before
training.

\begin{table}[t]
\caption{Dataset characteristics used by the experiments.}
\label{tab:data}
\centering
\begin{tabular}{lrrrl}
\toprule
Collection & Sections & Bins & Classes & Coverage \\
\midrule
GBM & 8 & 85,062 & 2 & 54.62--82.95\% \\
CAC & 8 & 1,481 & 3 & 89.98--100\% \\
\bottomrule
\end{tabular}
\end{table}

Each measured spectrum $s_i$ is total-ion-count (TIC) normalised,
\begin{equation}
x_i = \frac{s_i}{\sum_j s_{ij}+\epsilon},
\end{equation}
which reduces total-intensity scale differences while preserving relative bin
composition. The expert masks are retained separately for evaluation.

\subsection{Controlled VAE variants}

The centre-only control receives $x_i$. The uniform contextual model computes
the mean $c_i$ of the valid spectra in the eight-position Moore neighbourhood
and receives $[x_i,c_i]$. Missing neighbours contribute neither values nor
weights. The corrected-attention variant learns neighbour weights and applies
sqrt-bin input scaling before its scoring operation to avoid nearly uniform
scores caused by the small TIC-normalised magnitudes. All variants use a
512-unit hidden layer and a five-dimensional Gaussian latent representation.
The decoder reconstructs the central spectrum, so added context is useful only
if it helps explain the same target.

For a sample $x$, the objective is
\begin{equation}
\mathcal{L}=\mathcal{L}_{\mathrm{recon}}(x,\hat{x})+
\beta D_{\mathrm{KL}}\!\left(q_\phi(z\mid x,c)\|\mathcal{N}(0,I)\right).
\end{equation}
Matched comparisons use the same optimiser, schedule, split, epoch count and
seed. Production VAE runs use 100 epochs. Explicit adjacent-latent penalties
and Poisson augmentation were evaluated as development gates but were not
promoted: the former forced smoothness while degrading reconstruction, and the
latter did not reach its predeclared noisy-reconstruction improvement.

\subsection{Label-free nonlinear attribution}

After training, posterior means $\mu_i$ are frozen and a two-component GMM is
fitted without tissue labels. For pixel $i$, the scalar attribution target is
the posterior probability of its assigned component, $F_i=p(k_i\mid\mu_i)$.
For input $x$ and section-mean baseline $x'$, Integrated Gradients for bin $j$
is approximated by
\begin{equation}
\mathrm{IG}_j(x)=(x_j-x'_j)\frac{1}{M}
\sum_{m=1}^{M}\frac{\partial F(x'+\frac{m}{M}(x-x'))}{\partial x_j}.
\end{equation}
Absolute attributions are aggregated by $m/z$ across sampled pixels and GMM
components. Component-specific rankings are combined deterministically. For a
contextual model, contributions through the central and neighbourhood paths
are mapped to the same spectral bins before aggregation.

Completeness checks whether summed attributions approximate
$F(x)-F(x')$. Deletion faithfulness replaces highly ranked bins by baseline
values and measures the reduction in assigned-component posterior; a larger
drop indicates a more causally influential ranking. First-layer L2 importance
is retained as a deliberately simple linear comparator.

\subsection{Peak evaluation and statistical design}

For every $m/z$ bin, its ion image is correlated with each expert class mask
using Pearson's $r$. At a threshold $t\in\{0.3,0.4,0.5,0.6\}$, sufficiently
correlated bins form a reference-positive set. Precision and recall of a
selected peak list yield $F1_t$, and
\begin{equation}
\mathrm{mSCF1}=\frac{F1_{0.3}+F1_{0.4}+F1_{0.5}+F1_{0.6}}{4}.
\end{equation}
Within a section, compared methods select the same number of peaks. Results are
paired by section and exact two-sided Wilcoxon signed-rank tests are reported
where applicable. Section spread is not interpreted as patient-level
independence. Broad GBM and CAC validation uses seed 1; a targeted
GBM108-positive experiment repeats training with seeds 1--3 for centre-only,
uniform-mean and corrected-attention variants.
